\section{Related Work}
\label{sec:related_work}
\paragraph{Multi-Agent Systems.}
Recent LLM-based multi-agent systems improve task performance by decomposing problems across specialized or interacting agents. Early frameworks such as CAMEL~\citep{li2023camel} instantiate role-playing agents to study autonomous cooperation. MetaGPT~\citep{hong2023metagpt} and ChatDev~\citep{qian2024chatdev} further structure collaboration through software-engineering-inspired roles, communication protocols, and standardized workflows, demonstrating that specialization can reduce coordination complexity in multi-step tasks. AOrchestra~\citep{ruan2026aorchestra} dynamically creates task-specific sub-agents by instantiating their instructions, context, tools, and models; ToolOrchestra~\citep{su2025toolorchestra} trains a lightweight orchestrator to coordinate models and tools under accuracy, efficiency, and user-preference rewards; Squeeze Evolve~\citep{maheswaran2026squeeze} routes different stages of evolutionary inference to models of different cost and capability; and RecursiveMAS~\citep{yang2026recursive} replaces text-level communication with recursive latent-state transfer across agents.
Combee~\citep{li2026combee} also identifies centralized aggregation as a bottleneck under high parallelism, but focuses on scaling parallel prompt learning rather than decentralized agent coordination. Recent blackboard-style MAS~\citet{han2025exploring, salemi2025llm} use a shared board to let agents exchange messages or volunteer information, either for general reasoning or for data discovery.

These systems show that agent coordination can be scaled through dynamic delegation, learned orchestration, model routing, recursive collaboration, and shared blackboards. However, the orchestration-based methods primarily optimize how agents are created, routed, or composed, leaving the communication substrate itself unverified: intermediate outputs are passed, summarized, or aggregated without grounding each entry in its underlying evidence before downstream reuse; the blackboard-based methods write raw or lightly structured messages to the board and still route coordination through a selection step or a central poster. In contrast, \toolName{} treats the shared context as curated, verified state: agents coordinate fully asynchronously through a task queue with no central controller, every entry is admitted only after verification against its supporting evidence, and entries are stored as compact gists that unfold to finer detail on demand, keeping the shared state both reliable and scalable as the number of agents grows.

\paragraph{Programmatic Agentic Systems.}
A closely related class of systems equips language models with programmatic interfaces to external environments, enabling them to inspect state, invoke tools, manipulate artifacts, and recursively delegate computation. Recursive Language Models (RLMs)~\citep{zhang2025recursive} apply this paradigm to long-context reasoning by treating the prompt as an external environment and recursively selecting snippets for inspection. It performs this process through a code-mediated Read-Eval-Print Loop (REPL). Agentic coding systems such as Claude Code~\citep{anthropic_claude_code_2026} and Codex~\citep{openai_codex_2025} extend this idea to software environments, allowing models to read repositories, edit files, execute commands, and incorporate tool feedback. These approaches demonstrate that coupling language models with executable environments substantially enhances capability beyond single-pass generation. 

However, programmatic access alone does not solve the coordination problem. In many systems, planning is still routed through a main agent or local tool loop, and intermediate results are reused as raw observations or prose summaries without an admission-time grounding step. As a result, useful discoveries may remain local to one trajectory, while unsupported claims or softened constraints can propagate through later reasoning. \toolName{} is complementary to these systems: it does not replace tools, code execution, or REPL-style inspection, but provides a decentralized coordination layer in which tool-derived findings, failures, and partial solutions can be compressed, verified, and shared across agents. This complementarity is reflected in our OOLONG experiments (\S~\ref{sec:with_rlm}), where combining \toolName{} with RLM improves over either method alone.

\paragraph{Long-context LM systems.}
Another line of work addresses long-context reasoning by storing, compressing, or retrieving information beyond the model's current context window. \citet{mu2023learning} introduce a training-based approach that compresses long prompts into compact ``gist'' representations. Context-Folding~\citep{sun2025scaling} keeps long interactions within the context window by periodically compressing the conversation history. These methods reduce context length, but once fine-grained information is discarded, it cannot be reliably recovered when later reasoning requires precise details. ReadAgent~\citep{lee2024human} instead pairs compression with retrieval: it summarizes pages of text into short gist memories and looks up the original passages when a task needs details the gists omit. However, because lookup is triggered from lossy gists, ReadAgent can miss evidence whose relevance is not preserved in the compressed representation. 

Memory-augmented systems such as LongMem~\citep{wang2023augmenting}, MemGPT~\citep{packer2023memgpt}, Mem0~\citep{chhikara2025mem0}, and MemOS~\citep{li2025memos} introduce external memory mechanisms for LLMs to retain information across long interactions or task histories. For instance, MemGPT~\citep{packer2023memgpt} proposes an operating-system-inspired virtual context mechanism to decide what to keep in the active context and what to move into external memory. These systems make memory persistent, but retrieval is usually flat: a memory item is either selected into context or omitted. \toolName{} instead organizes shared information by abstraction level. Agents first reason over compact verified gists across the full problem, then selectively unfold relevant entries into detailed summaries and raw evidence. This hierarchy makes the shared context both global and recoverable.

